\centerline{\bf \Huge Домашнее задание.}
\centerline{\bf \Huge Вписанные четырёхугольники.}

\begin{flushright}
        \scriptsize{Окружность — душа геометрии. Познайте окружность,\\ и вы не только познаете геометрию, но и возвысите душу свою…\\И.Ф. Шарыгин}
\end{flushright}


Выпуклый четырёхугольник $ABCD$ можно вписать в окружность тогда и только тогда, когда выполнено хотя бы одно из эквивалентных условий:
\[
(1) \ \angle CAD = \angle CBD; \quad
(2) \ \angle ABD = \angle ACD; \quad
(3) \ \angle ABC + \angle ADC = 180^{\circ}.
\]

\begin{center}
\includegraphics[width=1\textwidth]{Летние Сборы/Сборы Элиста 2024/РВ/Geometry_9/Cyclic_qudrilateral.png}
\end{center}


\problem{1}
Шестиугольник $ABCDEF$ вписан в окружность. Докажите, что
\[\angle A + \angle C + \angle E = \angle B + \angle D + \angle F\]

\problem{2}
Шестиугольник \( A_1 A_2 A_3 A_4 A_5 A_6 \) вписан в окружность. Оказалось, что \( A_1 A_2 \parallel A_4 A_5 \) и \( A_2 A_3 \parallel A_5 A_6 \). Докажите, что \( A_3 A_4 \parallel A_6 A_1 \).

\problem{3}
Внутри треугольника $ABC$ отмечена точка $P$, а на сторонах $BC, CA, AB$ — точки $A_1, B_1, C_1$ соответственно. Докажите, что если два из трех четырехугольников $AB_1PC_1, BC_1PA_1, CA_1PB_1$ вписанные, то третий тоже вписанный.

\problem{$4^*$}
Окружность $\omega$ описана около остроугольного треугольника $ABC.$ Точка $H$ --- основание высоты $AH$ этого треугольника, а точка $G$ --- его точка пересечения медиан. Луч $GH$ пересекает $\omega$ в точке $P.$ Докажите, что $\angle BPH = \angle ABH$

\problem{$5^*$}
Точка $P$ выбирается на дуге $BC$ описанной около остроугольного треугольника ABC окружности, не содержащей точку $A$. Точка $Q$ внутри треугольника такова, что $\angle ABQ = \angle PBC$ и $\angle ACQ = \angle PCB.$ Докажите, что $AP = P Q$.